%%%%%%%%%%%%%%%%%%%%%%%%%%%%%%%%%%%%%%%%
% datoteka diploma-vzorec.tex
%
% vzorčna datoteka za pisanje diplomskega dela v formatu LaTeX
% na UL Fakulteti za računalništvo in informatiko
%
% vkup spravil Gašper Fijavž, december 2010
% 
%
%
% verzija 12. februar 2014 (besedilo teme, seznam kratic, popravki Gašper Fijavž)
% verzija 10. marec 2014 (redakcijski popravki Zoran Bosnić)
% verzija 11. marec 2014 (redakcijski popravki Gašper Fijavž)
% verzija 15. april 2014 (pdf/a 1b compliance, not really - just claiming, Damjan Cvetan, Gašper Fijavž)
% verzija 23. april 2014 (privzeto cc licenca)
% verzija 16. september 2014 (odmiki strain od roba)
% verzija 28. oktober 2014 (odstranil vpisno številko)
% verija 5. februar 2015 (Literatura v kazalu, online literatura)
% verzija 25. september 2015 (angl. naslov v izjavi o avtorstvu)
% verzija 26. februar 2016 (UL izjava o avtorstvu)
% verzija 16. april 2016 (odstranjena izjava o avtorstvu)
% verzija 5. junij 2016 (Franc Solina dodal vrstice, ki jih je označil s svojim imenom)


\documentclass[a4paper, 12pt]{book}
%\documentclass[a4paper, 12pt, draft]{book}  Nalogo preverite tudi z opcijo draft, ki vam bo pokazala, katere vrstice so predolge!



\usepackage[utf8x]{inputenc}   % omogoča uporabo slovenskih črk kodiranih v formatu UTF-8
\usepackage[slovene,english]{babel}    % naloži, med drugim, slovenske delilne vzorce
\usepackage[pdftex]{graphicx}  % omogoča vlaganje slik različnih formatov
\usepackage{fancyhdr}          % poskrbi, na primer, za glave strani
\usepackage{amssymb}           % dodatni simboli
\usepackage{amsmath}           % eqref, npr.
%\usepackage{hyperxmp}
\usepackage[hyphens]{url}  % dodal Solina
\usepackage{comment}       % dodal Solina

\usepackage[pdftex, colorlinks=true,
						citecolor=black, filecolor=black, 
						linkcolor=black, urlcolor=black,
						pagebackref=false, 
						pdfproducer={LaTeX}, pdfcreator={LaTeX}, hidelinks]{hyperref}

\usepackage{color}       % dodal Solina
\usepackage{soul}       % dodal Solina
\usepackage[numbers]{natbib}  % dodal Solina
\usepackage{tikz}
\usepackage{float}
\usepackage{longtable}

%%%%%%%%%%%%%%%%%%%%%%%%%%%%%%%%%%%%%%%%
%	DIPLOMA INFO
%%%%%%%%%%%%%%%%%%%%%%%%%%%%%%%%%%%%%%%%
\newcommand{\ttitle}{Analiza spletnih storitev Microsoft 365 in uporaba Microsoft Graph pri razvoju informacijskih sistemov}
\newcommand{\ttitleEn}{Analysis of Microsoft 365 web services and use of Microsoft Graph in information systems development}
\newcommand{\tsubject}{\ttitle}
\newcommand{\tsubjectEn}{\ttitleEn}
\newcommand{\tauthor}{Mia Erbus}
\newcommand{\tkeywords}{Microsoft 365, Microsoft Graph API, Microsoft Entra ID, ogrodje React}
\newcommand{\tkeywordsEn}{Microsoft 365, Microsoft Graph API, Microsoft Entra ID, React framework}


%%%%%%%%%%%%%%%%%%%%%%%%%%%%%%%%%%%%%%%%
%	HYPERREF SETUP
%%%%%%%%%%%%%%%%%%%%%%%%%%%%%%%%%%%%%%%%
\hypersetup{pdftitle={\ttitle}}
\hypersetup{pdfsubject=\ttitleEn}
\hypersetup{pdfauthor={\tauthor, me9223@student.uni-lj.si}}
\hypersetup{pdfkeywords=\tkeywordsEn}


 


%%%%%%%%%%%%%%%%%%%%%%%%%%%%%%%%%%%%%%%%
% postavitev strani
%%%%%%%%%%%%%%%%%%%%%%%%%%%%%%%%%%%%%%%%  

\addtolength{\marginparwidth}{-20pt} % robovi za tisk
\addtolength{\oddsidemargin}{40pt}
\addtolength{\evensidemargin}{-40pt}

\renewcommand{\baselinestretch}{1.3} % ustrezen razmik med vrsticami
\setlength{\headheight}{15pt}        % potreben prostor na vrhu
\renewcommand{\chaptermark}[1]%
{\markboth{\MakeUppercase{\thechapter.\ #1}}{}} \renewcommand{\sectionmark}[1]%
{\markright{\MakeUppercase{\thesection.\ #1}}} \renewcommand{\headrulewidth}{0.5pt} \renewcommand{\footrulewidth}{0pt}
\fancyhf{}
\fancyhead[LE,RO]{\sl \thepage} 
%\fancyhead[LO]{\sl \rightmark} \fancyhead[RE]{\sl \leftmark}
\fancyhead[RE]{\sc \tauthor}              % dodal Solina
\fancyhead[LO]{\sc Diplomska naloga}     % dodal Solina


\newcommand{\BibTeX}{{\sc Bib}\TeX}

%%%%%%%%%%%%%%%%%%%%%%%%%%%%%%%%%%%%%%%%
% naslovi
%%%%%%%%%%%%%%%%%%%%%%%%%%%%%%%%%%%%%%%%  


\newcommand{\autfont}{\Large}
\newcommand{\titfont}{\LARGE\bf}
\newcommand{\clearemptydoublepage}{\newpage{\pagestyle{empty}\cleardoublepage}}
\setcounter{tocdepth}{2}	      % globina kazala

%%%%%%%%%%%%%%%%%%%%%%%%%%%%%%%%%%%%%%%%
% konstrukti
%%%%%%%%%%%%%%%%%%%%%%%%%%%%%%%%%%%%%%%%  
\newtheorem{izrek}{Izrek}[chapter]
\newtheorem{trditev}{Trditev}[izrek]
\newenvironment{dokaz}{\emph{Dokaz.}\ }{\hspace{\fill}{$\Box$}}

%%%%%%%%%%%%%%%%%%%%%%%%%%%%%%%%%%%%%%%%%%%%%%%%%%%%%%%%%%%%%%%%%%%%%%%%%%%%%%%
%% PDF-A
%%%%%%%%%%%%%%%%%%%%%%%%%%%%%%%%%%%%%%%%%%%%%%%%%%%%%%%%%%%%%%%%%%%%%%%%%%%%%%%


%%%%%%%%%%%%%%%%%%%%%%%%%%%%%%%%%%%%%%%% 
% define medatata
%%%%%%%%%%%%%%%%%%%%%%%%%%%%%%%%%%%%%%%% 
\def\Title{\ttitle}
\def\Author{\tauthor, me9223@student.uni-lj.si}
\def\Subject{\ttitleEn}
\def\Keywords{\tkeywordsEn}

%%%%%%%%%%%%%%%%%%%%%%%%%%%%%%%%%%%%%%%% 
% \convertDate converts D:20080419103507+02'00' to 2008-04-19T10:35:07+02:00
%%%%%%%%%%%%%%%%%%%%%%%%%%%%%%%%%%%%%%%% 
\def\convertDate{%
    \getYear
}

{\catcode`\D=12
 \gdef\getYear D:#1#2#3#4{\edef\xYear{#1#2#3#4}\getMonth}
}
\def\getMonth#1#2{\edef\xMonth{#1#2}\getDay}
\def\getDay#1#2{\edef\xDay{#1#2}\getHour}
\def\getHour#1#2{\edef\xHour{#1#2}\getMin}
\def\getMin#1#2{\edef\xMin{#1#2}\getSec}
\def\getSec#1#2{\edef\xSec{#1#2}\getTZh}
\def\getTZh +#1#2{\edef\xTZh{#1#2}\getTZm}
\def\getTZm '#1#2'{%
    \edef\xTZm{#1#2}%
    \edef\convDate{\xYear-\xMonth-\xDay T\xHour:\xMin:\xSec+\xTZh:\xTZm}%
}

\expandafter\convertDate\pdfcreationdate 

%%%%%%%%%%%%%%%%%%%%%%%%%%%%%%%%%%%%%%%%
% get pdftex version string
%%%%%%%%%%%%%%%%%%%%%%%%%%%%%%%%%%%%%%%% 
\newcount\countA
\countA=\pdftexversion
\advance \countA by -100
\def\pdftexVersionStr{pdfTeX-1.\the\countA.\pdftexrevision}


%%%%%%%%%%%%%%%%%%%%%%%%%%%%%%%%%%%%%%%%
% XMP data
%%%%%%%%%%%%%%%%%%%%%%%%%%%%%%%%%%%%%%%%  
\usepackage{xmpincl}
\includexmp{pdfa-1b}

%%%%%%%%%%%%%%%%%%%%%%%%%%%%%%%%%%%%%%%%
% pdfInfo
%%%%%%%%%%%%%%%%%%%%%%%%%%%%%%%%%%%%%%%%  
\pdfinfo{%
    /Title    (\ttitle)
    /Author   (\tauthor, damjan@cvetan.si)
    /Subject  (\ttitleEn)
    /Keywords (\tkeywordsEn)
    /ModDate  (\pdfcreationdate)
    /Trapped  /False
}


%%%%%%%%%%%%%%%%%%%%%%%%%%%%%%%%%%%%%%%%%%%%%%%%%%%%%%%%%%%%%%%%%%%%%%%%%%%%%%%
%%%%%%%%%%%%%%%%%%%%%%%%%%%%%%%%%%%%%%%%%%%%%%%%%%%%%%%%%%%%%%%%%%%%%%%%%%%%%%%

\begin{document}
\selectlanguage{slovene}
\frontmatter
\setcounter{page}{1} %
\renewcommand{\thepage}{}       % preprecimo težave s številkami strani v kazalu
\newcommand{\sn}[1]{"`#1"'}                    % dodal Solina (slovenski narekovaji)

%%%%%%%%%%%%%%%%%%%%%%%%%%%%%%%%%%%%%%%%
%naslovnica
 \thispagestyle{empty}%
   \begin{center}
    {\large\sc Univerza v Ljubljani\\%
%      Fakulteta za elektrotehniko\\% za študijski program Multimedija
%      Fakulteta za upravo\\% za študijski program Upravna informatika
      Fakulteta za računalništvo in informatiko\\%
     Fakulteta za matematiko in fiziko\\% za študijski program Računalništvo in matematika
     }
    \vskip 10em%
    {\autfont \tauthor\par}%
    {\titfont \ttitle \par}%
    {\vskip 3em \textsc{DIPLOMSKO DELO\\[5mm]         % dodal Solina za ostale študijske programe
%    VISOKOŠOLSKI STROKOVNI ŠTUDIJSKI PROGRAM\\ PRVE STOPNJE\\ RAČUNALNIŠTVO IN INFORMATIKA}\par}%
%     UNIVERZITETNI  ŠTUDIJSKI PROGRAM\\ PRVE STOPNJE\\ RAČUNALNIŠTVO IN INFORMATIKA}\par}%
%    INTERDISCIPLINARNI UNIVERZITETNI\\ ŠTUDIJSKI PROGRAM PRVE STOPNJE\\ MULTIMEDIJA}\par}%
%    INTERDISCIPLINARNI UNIVERZITETNI\\ ŠTUDIJSKI PROGRAM PRVE STOPNJE\\ UPRAVNA INFORMATIKA}\par}%
   INTERDISCIPLINARNI UNIVERZITETNI\\ ŠTUDIJSKI PROGRAM PRVE STOPNJE\\ RAČUNALNIŠTVO IN MATEMATIKA}\par}%
    \vfill\null%
% izberite pravi habilitacijski naziv mentorja!
    {\large \textsc{Mentor}: doc. dr. Rok Rupnik\par}%
%   {\large \textsc{Somentor}:  viš. pred./doc./izr. prof./prof. dr.  Martin Krpan \par}%
    {\vskip 2em \large Ljubljana, 2026 \par}%
\end{center}
% prazna stran
%\clearemptydoublepage      % dodal Solina (izjava o licencah itd. se izpiše na hrbtni strani naslovnice)

%%%%%%%%%%%%%%%%%%%%%%%%%%%%%%%%%%%%%%%%
%copyright stran
\thispagestyle{empty}
\vspace*{8cm}

\noindent
{\sc Copyright}. 
Rezultati diplomske naloge so intelektualna lastnina avtorja in matične fakultete Univerze v Ljubljani.
Za objavo in koriščenje rezultatov diplomske naloge je potrebno pisno privoljenje avtorja, fakultete ter mentorja.

\begin{center}
\mbox{}\vfill
\emph{Besedilo je oblikovano z urejevalnikom besedil \LaTeX.}
\end{center}
% prazna stran
\clearemptydoublepage

%%%%%%%%%%%%%%%%%%%%%%%%%%%%%%%%%%%%%%%%
% stran 3 med uvodnimi listi
\thispagestyle{empty}
\
\vfill

\bigskip
\noindent\textbf{Kandidat:} Mia Erbus\\
\noindent\textbf{Naslov:} Analiza spletnih storitev Microsoft 365 in uporaba Microsoft Graph pri razvoju informacijskih sistemov\\
% vstavite ustrezen naziv študijskega programa!
\noindent\textbf{Vrsta naloge:} Diplomska naloga na univerzitetnem programu prve stopnje interdisciplinarnega študija Računalništvo in matematika \\
% izberite pravi habilitacijski naziv mentorja!
\noindent\textbf{Mentor:} doc. dr. Rok Rupnik\\
% \noindent\textbf{Somentor:} isto kot za mentorja

\bigskip
\noindent\textbf{Opis:}\\
Besedilo teme diplomskega dela študent prepiše iz študijskega informacijskega sistema, kamor ga je vnesel mentor. 
V nekaj stavkih bo opisal, kaj pričakuje od kandidatovega diplomskega dela. 
Kaj so cilji, kakšne metode naj uporabi, morda bo zapisal tudi ključno literaturo.

\bigskip
\noindent\textbf{Title:} Analysis of Microsoft 365 web services and use of Microsoft Graph in information systems development

\bigskip
\noindent\textbf{Description:}\\
opis diplome v angleščini

\vfill



\vspace{2cm}

% prazna stran
\clearemptydoublepage

% zahvala
\thispagestyle{empty}\mbox{}\vfill\null\it%
\noindent
Na tem mestu zapišite, komu se zahvaljujete za izdelavo diplomske naloge. Pazite, da ne boste koga pozabili. Utegnil vam bo zameriti. Temu se da izogniti tako, da celotno zahvalo izpustite.
\rm\normalfont

% prazna stran
\clearemptydoublepage

%%%%%%%%%%%%%%%%%%%%%%%%%%%%%%%%%%%%%%%%
% posvetilo, če sama zahvala ne zadošča :-)
% \thispagestyle{empty}\mbox{}{\vskip0.20\textheight}\mbox{}\hfill\begin{minipage}{0.55\textwidth}%
% Svoji dragi Alenčici.
% \normalfont\end{minipage}

% prazna stran
\clearemptydoublepage


%%%%%%%%%%%%%%%%%%%%%%%%%%%%%%%%%%%%%%%%
% kazalo
\pagestyle{empty}
\def\thepage{}% preprecimo tezave s stevilkami strani v kazalu
\tableofcontents{}


% prazna stran
\clearemptydoublepage

%%%%%%%%%%%%%%%%%%%%%%%%%%%%%%%%%%%%%%%%
% seznam kratic

\chapter*{Seznam uporabljenih kratic}  % spremenil Solina, da predolge vrstice ne gredo preko desnega roba

\begin{comment}
\begin{tabular}{l|l|l}
  {\bf kratica} & {\bf angleško} & {\bf slovensko} \\ \hline
  % after \\: \hline or \cline{col1-col2} \cline{col3-col4} ...
  {\bf CA} & classification accuracy & klasifikacijska točnost \\
  {\bf DBMS} & database management system & sistem za upravljanje podatkovnih baz \\
  {\bf SVM} & support vector machine & metoda podpornih vektorjev \\
  \dots & \dots & \dots \\
\end{tabular}
\end{comment}

\noindent\begin{tabular}{p{0.1\textwidth}|p{.4\textwidth}|p{.4\textwidth}}    % po potrebi razširi prvo kolono tabele na račun drugih dveh!
  {\bf kratica} & {\bf angleško}                             & {\bf slovensko} \\ \hline
  {\bf CA}      & classification accuracy               & klasifikacijska točnost \\
  {\bf DBMS} & database management system & sistem za upravljanje podatkovnih baz \\
  {\bf SVM}   & support vector machine              & metoda podpornih vektorjev \\
%  \dots & \dots & \dots \\
\end{tabular}


% prazna stran
\clearemptydoublepage

%%%%%%%%%%%%%%%%%%%%%%%%%%%%%%%%%%%%%%%%
% povzetek
\addcontentsline{toc}{chapter}{Povzetek}
\chapter*{Povzetek}

\noindent\textbf{Naslov:} \ttitle
\bigskip

\noindent\textbf{Avtor:} \tauthor
\bigskip

%\noindent\textbf{Povzetek:} 
\noindent V vzorcu je predstavljen postopek priprave diplomskega dela z uporabo okolja \LaTeX. Vaš povzetek mora sicer vsebovati približno 100 besed, ta tukaj je odločno prekratek.
Dober povzetek vključuje: (1) kratek opis obravnavanega problema, (2) kratek opis vašega pristopa za reševanje tega problema in (3) (najbolj uspešen) rezultat ali prispevek magistrske naloge.

\bigskip

\noindent\textbf{Ključne besede:} \tkeywords.
% prazna stran
\clearemptydoublepage

%%%%%%%%%%%%%%%%%%%%%%%%%%%%%%%%%%%%%%%%
% abstract
\selectlanguage{english}
\addcontentsline{toc}{chapter}{Abstract}
\chapter*{Abstract}

\noindent\textbf{Title:} \ttitleEn
\bigskip

\noindent\textbf{Author:} \tauthor
\bigskip

%\noindent\textbf{Abstract:} 
\noindent This sample document presents an approach to typesetting your BSc thesis using \LaTeX. 
A proper abstract should contain around 100 words which makes this one way too short.
\bigskip

\noindent\textbf{Keywords:} \tkeywordsEn.
\selectlanguage{slovene}
% prazna stran
\clearemptydoublepage

%%%%%%%%%%%%%%%%%%%%%%%%%%%%%%%%%%%%%%%%
\mainmatter
\setcounter{page}{1}
\pagestyle{fancy}

\chapter{Uvod}
Sodobna podjetja pri svojem poslovanju uporabljajo številne informacijske sisteme in komunikacijska orodja, kot so elektronska pošta, koledarji, sistemi za videokonference, dokumentni repozitoriji ter različne poslovne aplikacije.
Ker so te storitve pogosto med seboj ločene, se pojavljajo izzivi pri povezovanju podatkov, podvajanju funkcionalnosti in učinkovitem sodelovanju med uporabniki ter aplikacijami.

Razvoj oblačnega računalništva je omogočil prehod z lokalno vzdrževanih informacijskih sistemov na celovite oblačne platforme, ki združujejo različne poslovne storitve v enotno okolje.
Med najpomembnejšimi tovrstnimi reši-tvami je Microsoft 365, ki poleg pisarniških aplikacij vključuje tudi storitve za komunikacijo, sodelovanje, upravljanje dokumentov, videokonference ter upravljanje uporabniških identitet.

Ključni povezovalni element platforme Microsoft 365 predstavlja Microsoft Graph API, ki razvijalcem omogoča enoten programski dostop do različ-nih storitev in podatkov znotraj ekosistema Microsoft 365.
Preko enotnega vmesnika je mogoče dostopati do storitev, kot so Outlook, OneDrive, SharePoint, Teams in Entra ID, ter njihove funkcionalnosti vključevati v lastne informacijske rešitve. Microsoft Graph tako omogoča razvoj prilagojenih poslovnih aplikacij, integracijo podatkov med storitvami ter avtomatizacijo poslovnih procesov.

Kljub številnim prednostim razvoj aplikacij na osnovi Microsoft Graph API predstavlja tudi določene tehnične izzive. Razvijalci morajo razumeti mehanizme avtentikacije in avtorizacije, ki temeljijo na protokolu OAuth 2.0, upravljati dostopne žetone ter ustrezno konfigurirati model dovoljenj.
Poleg tega je treba upoštevati varnostne zahteve, omejitve števila zahtevkov ter obsežno in pogosto spreminjajočo se dokumentacijo.
Čeprav organizacije Microsoft 365 množično uporabljajo za elektronsko pošto, upravljanje dokumentov in spletno sodelovanje, ostajajo možnosti razvoja prilagojenih informacijskih sistemov ter avtomatizacije poslovnih procesov v številnih primerih še vedno premalo izkoriščene.

Motivacija za izdelavo diplomske naloge izhaja iz potrebe po poglobljenem razumevanju platforme Microsoft 365 kot sodobnega okolja za razvoj informacijskih sistemov ter preučevanju možnosti, ki jih za integracijo poslovnih procesov in storitev ponuja Microsoft Graph API.
Namen naloge je analizirati arhitekturo platforme Microsoft 365 in njen programski vmesnik Microsoft Graph, raziskati možnosti povezovanja poslovnih aplikacij ter ovrednotiti njegovo uporabnost pri razvoju informacijskih rešitev.
Teoretična spoznanja bodo podprta s praktično implementacijo demonstracijske aplikacije.

V okviru praktičnega dela bo razvita enostranska spletna aplikacija z uporabo ogrodja React, programskega jezika TypeScript ter razvojnega orodja Vite.
Za avtentikacijo uporabnikov in pridobivanje dostopnih žetonov bo uporabljena knjižnica Microsoft Authentication Library (MSAL).
Aplikacija bo komunicirala z vmesnikom Microsoft Graph API ter omogočala dostop do podatkov uporabniškega profila, pošiljanje elektronske pošte prek storitve Outlook, pregled vsebine uporabniškega prostora OneDrive, nalaganje dokumentov in ustvarjanje povezav za njihovo skupno rabo. Takšna implementacija bo omogočila praktično preverjanje zmogljivosti in omejitev platforme pri razvoju sodobnih spletnih aplikacij.

Diplomska naloga je razdeljena na več poglavij.
V drugem poglavju je predstavljena platforma Microsoft 365 in njene ključne storitve.
Tretje poglavje obravnava Microsoft Graph API, njegove osnovne značilnosti, strukturo zahtevkov, področja uporabe, varnostni mehanizmi, avtentikacija, avtorizacija in model dovoljenj.
Četrto poglavje opisuje načrtovanje in razvoj demonstracijske enostranske aplikacije.
V petem poglavju je predstavljena analiza rezultatov, prednosti in omejitve izbrane rešitve, ugotovitve naloge ter podaja možnosti za nadaljnje delo.


\chapter{Arhitektura platforme Microsoft 365}

Microsoft 365 je celovita oblačna platforma, ki temelji na večnajemniški (angl. multi-tenant) arhitekturi in uporabnikom omogoča dostop do aplikacij ter storitev prek interneta.
V takšnem modelu si več organizacij deli skupno programsko in infrastrukturno okolje, pri čemer so podatki posameznih organizacij logično ločeni in zaščiteni.
Večina storitev je uporabnikom na voljo v obliki modela SaaS (angl. Software as a Service), medtem ko platforma za svoje delovanje uporablja oblačno infrastrukturo Microsoft Azure.

Arhitektura Microsoft 365 je zasnovana kot skupek medsebojno povezanih storitev, ki zagotavljajo komunikacijo, sodelovanje, upravljanje dokumentov, identitetne storitve in poslovne integracije.
Ključne storitve, kot so Exchange Online, SharePoint Online, OneDrive for Business in Microsoft Teams, delujejo kot samostojne storitvene komponente, povezane prek skupnega identitetnega sistema in standardiziranih programskih vmesnikov.
Takšna zasnova omogoča visoko razširljivost, neodvisno posodabljanje posameznih storitev ter visoko razpoložljivost sistema.

Arhitekturo platforme Microsoft 365 lahko razdelimo na šest medsebojno povezanih slojev:

\begin{itemize}
\item uporabniški sloj,
\item aplikacijski sloj,
\item podatkovni sloj,
\item identitetni in varnostni sloj,
\item integracijski in razvojni sloj,
\item infrastrukturni sloj.
\end{itemize}


\section{Uporabniški sloj}

Uporabniški sloj predstavlja vstopno točko v okolje Microsoft 365 in vključuje vse odjemalske aplikacije, preko katerih uporabniki dostopajo do storitev platforme.
Dostop je mogoč preko spletnih brskalnikov, namiznih aplikacij, mobilnih aplikacij in integriranih odjemalcev operacijskega sistema Windows.

Med najpogosteje uporabljenimi aplikacijami so Microsoft Word, Excel, PowerPoint, Outlook in Teams.
Vse aplikacije uporabljajo enotno uporabniško identiteto, kar omogoča poenoteno prijavo in sinhronizacijo podatkov med napravami.
Komunikacija med odjemalci in storitvami praviloma poteka preko varnih protokolov HTTPS.

Pomembna značilnost uporabniškega sloja je podpora sodelovanju v realnem času.
Uporabniki lahko istočasno urejajo dokumente, sodelujejo na videokonferencah in uporabljajo skupne delovne prostore, pri čemer se spremembe sinhronizirajo preko oblačnih storitev platforme Microsoft 365.


\section{Aplikacijski sloj}

Aplikacijski sloj vsebuje poslovno logiko posameznih storitev Microsoft 365.
Med ključne storitve sodijo Exchange Online za elektronsko pošto in koledarje, SharePoint Online za upravljanje dokumentov in intranetne portale, OneDrive for Business za osebno hrambo datotek ter Microsoft Teams za komunikacijo in sodelovanje uporabnikov.

Dodatne storitve vključujejo Planner, To Do, Loop, Bookings, Viva Insights in Office Online. Čeprav posamezne storitve uporabljajo lastne podatkovne modele in funkcionalnosti, jih povezuje skupna identitetna infrastruktura ter Microsoft Graph API, ki omogoča enoten dostop do podatkov in funkcionalnosti znotraj celotnega ekosistema Microsoft 365.
Takšna arhitektura omogoča modularnost, razširljivost in visoko stopnjo integracije med posameznimi storitvami.


\section{Podatkovni sloj}

Podatkovni sloj je odgovoren za shranjevanje, obdelavo in zaščito podatkov uporabnikov.
Podatki so geografsko porazdeljeni med Microsoftove podatkovne centre in replicirani z namenom zagotavljanja visoke razpoložljivosti ter odpornosti proti izpadom storitev.
Platforma uporablja različne tipe podatkovnih shramb, prilagojene specifičnim potrebam posameznih storitev, kot so poštni predali, dokumenti, koledarski podatki in podatki o sodelovanju uporabnikov.

Pomembne funkcionalnosti podatkovnega sloja vključujejo upravljanje različic dokumentov, varnostno kopiranje podatkov, obnovo po izpadu ter šifriranje podatkov med prenosom in v mirovanju.~\cite{assurance}

\section{Identitetni in varnostni sloj}

Identitetni in varnostni sloj temelji na storitvi Microsoft Entra ID, ki zagotavlja centralizirano upravljanje identitet, avtentikacijo uporabnikov in nadzor dostopa do virov. Platforma podpira enotno prijavo (angl. Single Sign-On), večfaktorsko avtentikacijo (angl. Multi-Factor Authentication), upravljanje vlog in pravic ter pogojni dostop (angl. Conditional Access).~\cite{entraid}

Microsoft 365 uporablja sodobne varnostne mehanizme, ki vključujejo šifriranje komunikacije, zaščito pred izgubo podatkov (angl. Data Loss Prevention), revizijske dnevnike, zaznavanje varnostnih groženj in storitve Microsoft Defender.
Arhitektura temelji na načelih modela Zero Trust, ki predpostavlja, da noben uporabnik ali naprava nista samodejno zaupanja vredna, zato se vsak dostop preverja na podlagi identitete, lokacije, stanja naprave in varnostnih pravil organizacije.~\cite{zerotrust}


\section{Integracijski in razvojni sloj}

Integracijski in razvojni sloj omogoča povezovanje storitev Microsoft 365 z zunanjimi informacijskimi sistemi ter razvoj lastnih aplikacij.
Osrednjo vlogo pri tem ima Microsoft Graph, ki predstavlja enotno programsko vstopno točko za dostop do podatkov in storitev znotraj okolja Microsoft 365.~\cite{graphoverview}

Preko Microsoft Graph API je mogoče dostopati do uporabniških profilov, elektronske pošte, koledarjev, datotek, skupin, nalog, komunikacije v Teams in drugih organizacijskih podatkov.
Avtentikacija aplikacij temelji na standardih OAuth 2.0 in OpenID Connect, razvijalcem pa so na voljo programski kompleti SDK za različne programske jezike.
Takšna zasnova omogoča razvoj prilagojenih poslovnih rešitev, avtomatizacijo procesov in integracijo Microsoft 365 z obstoječimi informacijskimi sistemi.~\cite{graphapi}


\section{Infrastrukturni sloj}

Infrastrukturni sloj temelji na globalni infrastrukturi platforme Microsoft Azure, ki zagotavlja računske vire, omrežne storitve, podatkovne centre in mehanizme za zagotavljanje razpoložljivosti storitev. Microsoft 365 deluje v globalno porazdeljenem okolju, ki omogoča geografsko redundanco podatkov, samodejno obnovo ob izpadih ter prilagajanje zmogljivosti glede na obremenitev sistema.~\cite{assurance}

Takšna arhitektura omogoča zanesljivo delovanje storitev za milijone uporabnikov po vsem svetu ter predstavlja temelj za razširljivost, varnost in neprekinjeno razpoložljivost platforme Microsoft 365.


\begin{figure}[H]
\centering

\begin{tikzpicture}[
    layer/.style={
        rectangle,
        draw=black,
        rounded corners=3pt,
        minimum width=13cm,
        minimum height=1.2cm,
        align=center,
        fill=gray!10
    },
    font=\small
]

\node[layer] (user) at (0,0)
{
\textbf{UPORABNIŠKI SLOJ}\\
Outlook, Teams, Word, Excel, spletne in mobilne aplikacije
};

\node[layer] (app) at (0,-1.8)
{
\textbf{APLIKACIJSKI SLOJ}\\
Microsoft 365 aplikacije in storitve
};

\node[layer] (data) at (0,-3.6)
{
\textbf{PODATKOVNI SLOJ}\\
OneDrive, SharePoint, Exchange Online, podatkovne storitve
};

\node[layer] (identity) at (0,-5.4)
{
\textbf{IDENTITETNI IN VARNOSTNI SLOJ}\\
Microsoft Entra ID, MFA, Conditional Access, varnostne politike
};

\node[layer] (integration) at (0,-7.2)
{
\textbf{INTEGRACIJSKI IN RAZVOJNI SLOJ}\\
Microsoft Graph API, SDK, Power Platform
};

\node[layer] (infra) at (0,-9)
{
\textbf{INFRASTRUKTURNI SLOJ}\\
Microsoft Azure, podatkovni centri, omrežna infrastruktura
};

\draw[->, thick] (user) -- (app);
\draw[->, thick] (app) -- (data);
\draw[->, thick] (data) -- (identity);
\draw[->, thick] (identity) -- (integration);
\draw[->, thick] (integration) -- (infra);

\end{tikzpicture}

\caption{Večslojna arhitektura platforme Microsoft 365}
\label{fig:m365-arhitektura}

\end{figure}

\chapter{Microsoft Graph}

Microsoft Graph predstavlja osrednji programski vmesnik (API) ekosistema Microsoft 365 in identitetne platforme Microsoft Entra.
Gre za enoten programski model, ki razvijalcem omogoča dostop do podatkov, storitev in funkcionalnosti različnih Microsoftovih oblačnih rešitev prek skupne vstopne točke.
Namesto uporabe ločenih programskih vmesnikov za posamezne storitve, kot so Outlook, SharePoint, OneDrive ali Microsoft Teams, Microsoft Graph zagotavlja poenoten način dostopa do virov z uporabo skupnega identitetnega sistema, enotnega modela dovoljenj ter standardiziranih podatkovnih struktur.

\section{Osnovne značilnosti Microsoft Graph}

Microsoft Graph deluje kot posredni sloj med odjemalskimi aplikacijami in Microsoftovimi oblačnimi storitvami.
Aplikacije preko vmesnika pošiljajo zahteve za dostop do podatkov in funkcionalnosti posameznih storitev, medtem ko Microsoft Graph skrbi za njihovo usmerjanje, preverjanje pravic dostopa ter vračanje rezultatov v standardizirani obliki.
Takšen pristop poenostavlja razvoj aplikacij, zmanjšuje kompleksnost integracij in omogoča enotno upravljanje dostopa do različnih storitev znotraj ekosistema Microsoft 365.

Vsi zahtevki se izvajajo preko centralne vstopne točke:

\begin{verbatim}
https://graph.microsoft.com
\end{verbatim}

Microsoft Graph temelji na arhitekturnem slogu REST (angl. Representational State Transfer) in za komunikacijo uporablja protokol HTTP.
Pri dostopu do virov uporablja standardne metode HTTP, med katerimi so najpogosteje uporabljene metoda GET za pridobivanje podatkov, POST za ustvarjanje novih virov, PATCH za posodabljanje obstoječih virov ter DELETE za njihovo odstranjevanje.
Takšna zasnova omogoča enostavno uporabo v različnih programskih okoljih in platformah.

Podatki se med odjemalcem in strežnikom izmenjujejo v formatu JSON (angl. JavaScript Object Notation), ki predstavlja de facto standard za izmenjavo podatkov v sodobnih spletnih storitvah.
Zaradi svoje preprostosti, berljivosti in široke podpore omogoča učinkovito integracijo Microsoft Graph API z namiznimi, spletnimi, mobilnimi in strežniškimi aplikacijami.

Pomembna prednost Microsoft Graph je tudi uporaba enotnega podatkovnega modela, ki omogoča dosledno predstavitev uporabnikov, skupin, dokumentov, sporočil, koledarskih dogodkov in drugih poslovnih objektov ne glede na izvorno storitev.
S tem Microsoft Graph predstavlja temeljni integracijski mehanizem za razvoj sodobnih aplikacij na platformi Microsoft 365.

\section{Arhitektura Microsoft Graph}

Arhitektura Microsoft Graph temelji na več medsebojno povezanih komponentah, ki skupaj omogočajo varen in poenoten dostop do podatkov ter storitev znotraj ekosistema Microsoft 365.
Med ključne gradnike arhitekture sodijo odjemalske aplikacije, identitetna platforma Microsoft Entra ID, storitev Microsoft Graph kot osrednja vstopna točka za dostop do virov ter posamezne storitve Microsoft 365, ki zagotavljajo poslovne funkcionalnosti in podatke.

Odjemalska aplikacija predstavlja izhodiščno točko komunikacije z Microsoft Graph API. Gre lahko za spletno aplikacijo, enostransko spletno aplikacijo (angl. Single Page Application – SPA), mobilno aplikacijo, namizno aplikacijo ali strežniško rešitev. Pred dostopom do zaščitenih virov mora aplikacija pridobiti ustrezno identiteto in dovoljenja za izvajanje zahtevkov.

Proces avtentikacije in avtorizacije poteka preko storitve Microsoft Entra ID, ki uporablja industrijska standarda OAuth 2.0 in OpenID Connect. Po uspešni prijavi uporabnika ali aplikacije Entra ID izda dostopni žeton (angl. access token), ki vsebuje informacije o identiteti prijavljenega uporabnika, dodeljenih dovoljenjih ter obsegu dostopa do posameznih virov.
Microsoft Graph uporablja te informacije za preverjanje avtorizacije in nadzor dostopa do podatkov.

Po pridobitvi dostopnega žetona aplikacija pošilja HTTP-zahtevke na Microsoft Graph.
Ta deluje kot centralizirana vstopna točka oziroma prehodni sloj (angl. gateway), ki sprejema zahtevke, preverja njihovo veljavnost ter jih usmerja do ustreznih storitev znotraj Microsoftovega oblačnega okolja.
Med najpomembnejše storitve, dostopne preko Microsoft Graph, sodijo Exchange Online, SharePoint Online, OneDrive for Business, Microsoft Teams, Planner, Outlook Calendar in Microsoft Entra ID.

Takšna arhitektura razvijalcem omogoča uporabo enotnega programskega vmesnika za dostop do različnih storitev in podatkovnih virov, ne glede na njihovo dejansko lokacijo ali implementacijo.
S tem se bistveno zmanjša kompleksnost razvoja integracijskih rešitev, saj aplikacijam ni treba komunicirati z več ločenimi programskimi vmesniki.
Poleg tega enoten model avtentikacije, dovoljenj in podatkovnih objektov poenostavlja razvoj, vzdrževanje in nadzor nad aplikacijami, ki uporabljajo storitve Microsoft 365.

Pomemben del arhitekture predstavljata tudi podatkovni in varnostni sloj, ki zagotavljata zaščito uporabniških podatkov, šifriranje komunikacije, revizijske sledi ter skladnost z varnostnimi in regulatornimi zahtevami.
Microsoft Graph tako ne deluje zgolj kot komunikacijski vmesnik, temveč kot osrednji integracijski mehanizem, ki povezuje identitetne storitve, poslovne aplikacije in podatkovne vire znotraj celotnega Microsoftovega oblačnega ekosistema.


\section{Avtentikacija in avtorizacija}

Dostop do virov Microsoft Graph je zaščiten z varnostnim modelom, ki temelji na standardih OAuth 2.0 in OpenID Connect.
Pred uporabo Microsoft Graph mora biti aplikacija registrirana v identitetni platformi Microsoft Entra ID, kjer pridobi enolične identifikatorje in poverilnice, potrebne za avtentikacijo in avtorizacijo. Med najpomembnejše identifikacijske podatke sodijo identifikator aplikacije (angl. Application (Client) ID), identifikator imenika (angl. Directory (Tenant) ID) ter po potrebi skrivni ključ aplikacije (angl. Client Secret) ali digitalno potrdilo.

Postopek dostopa do zaščitenih virov se začne z avtentikacijo uporabnika ali aplikacije pri storitvi Microsoft Entra ID.
Po uspešno izvedenem postopku avtentikacije Entra ID izda dostopni žeton (angl. access token), ki je običajno zapisan v obliki JSON Web Token (angl. JWT).
Žeton vsebuje informacije o identiteti uporabnika ali aplikacije, dodeljenih dovoljenjih ter obdobju veljavnosti.
Pri komunikaciji z Microsoft Graph se dostopni žeton posreduje v HTTP glavi Authorization:

\begin{verbatim}
Authorization: Bearer <access_token>
\end{verbatim}

Ob prejemu zahtevka Microsoft Graph preveri veljavnost dostopnega žetona, identiteto uporabnika ali aplikacije, dodeljena dovoljenja ter morebitne varnostne politike organizacije.
Šele po uspešno izvedenem postopku preverjanja je aplikaciji omogočen dostop do zahtevanih virov.

Model avtorizacije temelji na sistemu dovoljenj (angl. permissions oziroma scopes), ki določajo obseg dostopa do posameznih storitev in podatkov. Microsoft Graph podpira dve osnovni vrsti dovoljenj:

\begin{itemize}
\item delegirana dovoljenja (angl. Delegated Permissions),
\item aplikacijska dovoljenja (angl. Application Permissions).
\end{itemize}

Delegirana dovoljenja se uporabljajo v primerih, ko aplikacija deluje v imenu prijavljenega uporabnika.
Obseg dostopa je v tem primeru omejen z dovoljenji uporabnika in dovoljenji, ki so bila dodeljena aplikaciji.
Aplikacija lahko tako dostopa le do virov, do katerih ima dostop tudi prijavljeni uporabnik.

Aplikacijska dovoljenja omogočajo izvajanje operacij brez prisotnosti prijavljenega uporabnika. Takšen način dostopa se najpogosteje uporablja pri strežniških procesih, avtomatizaciji opravil, sinhronizaciji podatkov in drugih ozadnih storitvah. Ker aplikacija v tem primeru dostopa do podatkov neposredno, so aplikacijska dovoljenja praviloma strožje nadzorovana in zahtevajo administratorsko soglasje.

Primeri pogosto uporabljenih dovoljenj v Microsoft Graph so:

\begin{itemize}
\item User.Read – dostop do osnovnih podatkov prijavljenega uporabnika,
\item Mail.Send – pošiljanje elektronske pošte v imenu uporabnika,
\item Calendars.Read – branje koledarskih dogodkov,
\item Files.ReadWrite – branje in spreminjanje datotek,
\item Sites.Read.All – branje podatkov spletnih mest SharePoint.
\end{itemize}

Takšen model omogoča natančno določanje pravic dostopa in predstavlja pomemben mehanizem za zagotavljanje varnosti ter zaščite podatkov znotraj okolja Microsoft 365.


\section{Model soglasja}

Microsoft Graph uporablja model soglasja (angl. Consent Framework), ki omogoča nadzor nad dostopom aplikacij do podatkov in storitev znotraj okolja Microsoft 365. Preden lahko aplikacija uporablja določena dovoljenja, mora uporabnik ali skrbnik organizacije izrecno potrditi, da aplikaciji zaupa in ji dovoljuje dostop do zahtevanih virov.

Ob prvi prijavi aplikacija uporabniku prikaže seznam dovoljenj, ki jih potrebuje za svoje delovanje. Uporabnik lahko ta dovoljenja potrdi ali zavrne. Postopek potrditve dovoljenj se imenuje uporabniško soglasje (angl. user consent). Po potrditvi Microsoft Entra ID aplikaciji izda dostopni žeton, ki vsebuje odobrena dovoljenja.

Pri dovoljenjih z višjo stopnjo privilegijev uporabniško soglasje ni dovolj. V takšnih primerih mora dovoljenja odobriti skrbnik organizacije, kar imenujemo administratorsko soglasje (angl. admin consent). Administratorsko soglasje je običajno potrebno za aplikacijska dovoljenja ter za delegirana dovoljenja, ki omogočajo dostop do podatkov drugih uporabnikov ali organizacijskih virov.

Microsoft Entra ID omogoča centralizirano upravljanje soglasij in dovoljenj preko administrativnega portala. Skrbniki lahko pregledujejo, katerim aplikacijam so bila dodeljena dovoljenja, jih po potrebi prekličejo ter določajo politike, ki omejujejo možnost uporabniškega soglasja. Takšen pristop zmanjšuje tveganje nepooblaščenega dostopa do podatkov in omogoča učinkovitejše upravljanje varnosti v organizaciji.

Proces dodeljevanja dovoljenj v okolju Microsoft Graph vključuje naslednje korake:

\begin{enumerate}
\item registracijo aplikacije v Microsoft Entra ID,
\item konfiguracijo zahtevanih dovoljenj,
\item pridobitev soglasja uporabnika ali skrbnika,
\item pridobitev dostopnega žetona,
dostop do virov preko Microsoft Graph API.
\end{enumerate}

Pomembna značilnost modela soglasja je načelo najmanjših privilegijev (angl. Principle of Least Privilege), ki priporoča, da aplikacija zahteva le tista dovoljenja, ki so nujno potrebna za njeno delovanje.
Takšen pristop zmanjšuje varnostna tveganja in omejuje morebitne posledice zlorabe uporabniških računov ali aplikacijskih poverilnic.

V poslovnih okoljih je model soglasja še posebej pomemben, saj organizacijam omogoča nadzor nad aplikacijami tretjih ponudnikov ter zagotavlja skladnost z internimi varnostnimi politikami in regulatornimi zahtevami.
Zaradi tega predstavlja pomemben element celotnega varnostnega modela platforme Microsoft 365 in Microsoft Graph.

\section{Endpointi Microsoft Graph}
Microsoft Graph omogoča dostop do podatkov in funkcionalnosti različnih storitev Microsoft 365 preko enotne strukture URL-naslovov oziroma endpointov. Vsak endpoint predstavlja določen vir ali zbirko virov, kot so uporabniki, skupine, elektronska pošta, koledarji in datoteke.

Splošna oblika URL-naslova v Microsoft Graph je:

\begin{verbatim}
https://graph.microsoft.com/{version}/{resource}
\end{verbatim}
kjer parameter \textit{version} določa različico API-ja (v1.0 ali beta), parameter \textit{resource} pa predstavlja zahtevani vir.

Med najpogosteje uporabljene vire v Microsoft Graph sodijo:

\begin{itemize}
    \item \texttt{/users} -- upravljanje uporabnikov in organizacijskih podatkov;
    \item \texttt{/groups} -- upravljanje skupin in članstev;
    \item \texttt{/messages} -- dostop do elektronske pošte;
    \item \texttt{/events} -- upravljanje koledarjev in dogodkov;
    \item \texttt{/drive} -- dostop do datotek v OneDrive;
    \item \texttt{/sites} -- dostop do vsebin SharePoint;
    \item \texttt{/teams} -- upravljanje ekip in komunikacije v Teams.
\end{itemize}

Primer pridobivanja podatkov o trenutno prijavljenem uporabniku:

\begin{verbatim}
GET https://graph.microsoft.com/v1.0/me
\end{verbatim}


Microsoft Graph podpira tudi uporabo poizvedbenih parametrov OData, ki omogočajo filtriranje, razvrščanje in omejevanje rezultatov. Med najpogosteje uporabljenimi parametri so \texttt{\$select}, \texttt{\$filter}, \texttt{\$expand}, \texttt{\$orderby} in \texttt{\$top}.


Primer pridobivanja izbranih atributov uporabnika:

\begin{verbatim}
GET https://graph.microsoft.com/v1.0/me?$select=id,displayName,mail
\end{verbatim}


Enotna struktura endpointov in skupni podatkovni model predstavljata eno izmed ključnih prednosti Microsoft Graph, saj razvijalcem omogočata dostop do različnih storitev Microsoft 365 preko enotnega programskega vmesnika ter poenostavljata razvoj integracijskih rešitev.


\section{Ključne storitve Microsoft Graph}

Microsoft Graph omogoča dostop do podatkov in funkcionalnosti številnih storitev znotraj ekosistema Microsoft 365.
Preko enotnega programskega vmesnika lahko aplikacije uporabljajo identitetne storitve, elektronsko pošto, koledarje, dokumentne repozitorije, komunikacijska orodja in druge poslovne vire.
Takšen pristop poenostavlja razvoj informacijskih sistemov, saj razvijalcem omogoča uporabo enotnega modela avtentikacije, avtorizacije in dostopa do podatkov ne glede na storitev, s katero komunicirajo.

Ena izmed ključnih prednosti Microsoft Graph je možnost dostopa do različnih storitev Microsoft 365 preko enotnega programskega vmesnika in skupnega podatkovnega modela. Namesto uporabe ločenih API-jev za posamezne storitve lahko aplikacije uporabljajo enotne REST končne točke, skupen sistem avtorizacije na osnovi OAuth 2.0 ter standardizirane podatkovne objekte. Takšna arhitektura poenostavi razvoj integracij, zmanjšuje kompleksnost upravljanja dostopov in omogoča učinkovito povezovanje podatkov med različnimi storitvami znotraj ekosistema Microsoft 365.

V nadaljevanju so predstavljene najpomembnejše storitve in endpointi, ki so dostopni preko Microsoft Graph, ter njihove značilnosti z vidika razvoja poslovnih aplikacij in informacijskih sistemov.

\subsection{Uporabniki}

Vir \texttt{User} predstavlja osrednji objekt identitetnega sistema Microsoft Entra ID in omogoča dostop do podatkov o uporabnikih organizacije. Vsak uporabnik je predstavljen kot samostojen objekt z enoličnim identifikatorjem (\texttt{id}), ki ostane nespremenjen skozi celoten življenjski cikel računa.

Microsoft Graph omogoča pridobivanje osnovnih atributov uporabnika, kot so ime, elektronski naslov, delovno mesto, oddelek, telefonske številke in lokacija. Poleg osnovnih podatkov so dostopni tudi organizacijski odnosi, kot sta nadrejeni uporabnik (\texttt{manager}) in seznam podrejenih zaposlenih (\texttt{directReports}), kar omogoča gradnjo organizacijskih diagramov in sistemov za upravljanje kadrovskih procesov.

Pomembna funkcionalnost je dostop do uporabniških virov preko relacijskih povezav. Iz uporabniškega objekta je mogoče neposredno dostopati do poštnega predala, koledarjev, datotek v OneDrive, članstev v skupinah, ekip Teams in drugih povezanih storitev. Primer takšne povezave je končna točka: \texttt{/users/{id}/drive}, ki omogoča dostop do uporabnikovega prostora OneDrive.

Microsoft Graph podpira tudi upravljanje življenjskega cikla uporabnikov. Administratorji lahko ustvarjajo nove račune, spreminjajo njihove lastnosti, dodeljujejo licence, upravljajo članstva v skupinah ter izvajajo postopke obnovitve ali trajnega brisanja računov. Za večje organizacije je pomembna tudi podpora za množične operacije in sinhronizacijo identitet z lokalnimi imeniki.

\subsection{Skupine}

Vir \texttt{Group} predstavlja logično združevanje uporabnikov in naprav ter je eden najpomembnejših mehanizmov za upravljanje dostopa v okolju Microsoft 365. Microsoft Graph podpira več vrst skupin, med katerimi sta najpogosteje uporabljeni Microsoft 365 skupina in varnostna skupina (angl. Security Group).

Microsoft 365 skupina poleg članstva samodejno zagotavlja tudi skupni poštni predal, SharePoint mesto, koledar, OneNote beležnico in druge povezane vire. Zaradi tega predstavlja temelj sodelovalnih okolij v storitvah Teams, Planner in Outlook.

Microsoft Graph omogoča upravljanje članov skupin preko relacijskih virov:

\begin{verbatim}
/groups/{id}/members
/groups/{id}/owners
\end{verbatim}

Podprto je dodajanje in odstranjevanje članov, pridobivanje informacij o lastnikih skupine ter upravljanje dinamičnega članstva. Dinamične skupine uporabljajo pravila na osnovi atributov uporabnikov, kar omogoča avtomatsko vključevanje ali izključevanje članov brez ročnega posredovanja administratorjev.

Pri večjih sistemih je pomembna tudi možnost pridobivanja transitivega članstva, kjer Microsoft Graph vrne vse skupine, v katere je uporabnik vključen neposredno ali posredno preko ugnezdenih skupin.

\begin{longtable}{|p{7cm}|p{7.5cm}|}
\caption{Microsoft Graph API endpointi za delo s skupinami}
\label{tab:group-endpoints} \\
\hline
\textbf{Endpoint} & \textbf{Opis} \\
\hline
\endfirsthead

\multicolumn{2}{c}%
{{\tablename\ \thetable{} -- nadaljevanje}} \\
\hline
\textbf{Endpoint} & \textbf{Opis} \\
\hline
\endhead

\hline
\multicolumn{2}{r}{{Nadaljevanje na naslednji strani}} \\
\endfoot

\hline
\endlastfoot

\multicolumn{2}{|l|}{\textbf{Skupine}} \\
\hline

\texttt{GET /groups} &
Pridobi seznam skupin v organizaciji. \\
\hline

\texttt{POST /groups} &
Ustvari novo skupino. \\
\hline

\texttt{GET /groups/\{group-id\}} &
Pridobi podatke o določeni skupini. \\
\hline

\texttt{PATCH /groups/\{group-id\}} &
Posodobi lastnosti skupine. \\
\hline

\texttt{DELETE /groups/\{group-id\}} &
Izbriše skupino. \\
\hline

\multicolumn{2}{|l|}{\textbf{Člani skupine}} \\
\hline

\texttt{GET /groups/\{group-id\}/members} &
Pridobi člane skupine. \\
\hline

\texttt{POST /groups/\{group-id\}/members
/\$ref} &
Doda člana v skupino. \\
\hline

\texttt{DELETE /groups/\{group-id\}/members
/\{directory-object-id\}/\$ref} &
Odstrani člana iz skupine. \\
\hline

\multicolumn{2}{|l|}{\textbf{Lastniki skupine}} \\
\hline

\texttt{GET /groups/\{group-id\}/owners} &
Pridobi lastnike skupine. \\
\hline

\texttt{POST /groups/\{group-id\}/owners
/\$ref} &
Doda lastnika skupini. \\
\hline

\texttt{DELETE /groups/\{group-id\}/owners
/\{directory-object-id\}/\$ref} &
Odstrani lastnika skupine. \\
\hline

\multicolumn{2}{|l|}{\textbf{Pogovori}} \\
\hline

\texttt{GET /groups/\{group-id\}
/conversations} &
Pridobi pogovore skupine. \\
\hline

\texttt{POST /groups/\{group-id\}
/conversations} &
Ustvari nov pogovor. \\
\hline

\texttt{GET /groups/\{group-id\}
/conversations/\{conversation-id\}} &
Pridobi določen pogovor. \\
\hline

\texttt{DELETE /groups/\{group-id\}
/conversations/\{conversation-id\}} &
Izbriše pogovor. \\
\hline

\multicolumn{2}{|l|}{\textbf{Niti pogovorov}} \\
\hline

\texttt{GET /groups/\{group-id\}/threads} &
Pridobi niti pogovorov skupine. \\
\hline

\texttt{GET ../threads/\{thread-id\}} &
Pridobi določeno nit. \\
\hline

\texttt{POST ../threads/\{thread-id\}/reply} &
Doda odgovor v nit pogovora. \\
\hline

\multicolumn{2}{|l|}{\textbf{Objave}} \\
\hline

\texttt{GET /groups/\{group-id\}/threads
/\{thread-id\}/posts} &
Pridobi objave v niti. \\
\hline

\texttt{GET /groups/\{group-id\}/threads
/\{thread-id\}/posts/\{post-id\}} &
Pridobi določeno objavo. \\
\hline

\multicolumn{2}{|l|}{\textbf{Datoteke}} \\
\hline

\texttt{GET /groups/\{group-id\}/drive} &
Pridobi dokumentno knjižnico skupine. \\
\hline

\texttt{GET ../drive/root} &
Pridobi korensko mapo dokumentov skupine. \\
\hline

\texttt{GET ../drive/items/\{item-id\}} &
Pridobi datoteko ali mapo skupine. \\
\hline

\multicolumn{2}{|l|}{\textbf{Koledarji}} \\
\hline

\texttt{GET /groups/\{group-id\}/calendar} &
Pridobi privzeti koledar skupine. \\
\hline

\texttt{GET /groups/\{group-id\}/events} &
Pridobi dogodke skupine. \\
\hline

\texttt{POST /groups/\{group-id\}/events} &
Ustvari nov dogodek skupine. \\
\hline

\texttt{GET ../events/\{event-id\}} &
Pridobi določen dogodek. \\
\hline

\texttt{PATCH ../events/\{event-id\}} &
Posodobi dogodek. \\
\hline

\texttt{DELETE ../events/\{event-id\}} &
Izbriše dogodek. \\
\hline

\multicolumn{2}{|l|}{\textbf{Zvezki OneNote}} \\
\hline

\texttt{GET /groups/\{group-id\}/onenote
/notebooks} &
Pridobi zvezke OneNote skupine. \\
\hline

\texttt{POST /groups/\{group-id\}/onenote
/notebooks} &
Ustvari nov zvezek OneNote. \\
\hline

\multicolumn{2}{|l|}{\textbf{Spletna mesta SharePoint}} \\
\hline

\texttt{GET /groups/\{group-id\}/sites/root} &
Pridobi povezano SharePoint spletno mesto. \\
\hline

\multicolumn{2}{|l|}{\textbf{Ekipa Teams}} \\
\hline

\texttt{GET /groups/\{group-id\}/team} &
Pridobi ekipo Teams, povezano s skupino. \\
\hline

\texttt{PUT /groups/\{group-id\}/team} &
Ustvari ekipo Teams iz skupine. \\
\hline

\multicolumn{2}{|l|}{\textbf{Nastavitve in življenjski cikel}} \\
\hline

\texttt{POST /groups/\{group-id\}/renew} &
Podaljša življenjsko dobo skupine. \\
\hline

\texttt{GET /groups/\{group-id\}/settings} &
Pridobi nastavitve skupine. \\
\hline

\texttt{GET /groups/\{group-id\}
/transitiveMembers} &
Pridobi vse neposredne in posredne člane skupine. \\
\hline

\texttt{GET /groups/\{group-id\}
/transitiveMemberOf} &
Pridobi vse skupine, katerih član je skupina. \\
\hline

\end{longtable}

\subsection{Elektronska pošta}

\begin{longtable}{|p{7cm}|p{7.5cm}|}
\caption{Microsoft Graph API endpointi za delo z elektronsko pošto}
\label{tab:mail-endpoints}\\
\hline
\textbf{Endpoint} & \textbf{Opis} \\
\hline
\endfirsthead

\hline
\textbf{Endpoint} & \textbf{Opis} \\
\hline
\endhead

\hline
\endfoot

\hline
\endlastfoot

\multicolumn{2}{|l|}{\textbf{Sporočila}} \\
\hline

\texttt{GET /me/messages} &
Pridobi seznam sporočil prijavljenega uporabnika. \\
\hline

\texttt{POST /me/messages} &
Ustvari novo sporočilo v obliki osnutka. \\
\hline

\texttt{GET /me/messages/\{message-id\}} &
Pridobi podrobnosti določenega sporočila. \\
\hline

\texttt{PATCH /me/messages/\{message-id\}} &
Posodobi lastnosti sporočila. \\
\hline

\texttt{DELETE /me/messages/\{message-id\}} &
Izbriše sporočilo. \\
\hline

\multicolumn{2}{|l|}{\textbf{Pošiljanje sporočil}} \\
\hline

\texttt{POST /me/sendMail} &
Pošlje novo elektronsko sporočilo. \\
\hline

\texttt{POST /me/messages/\{message-id\}
/send} &
Pošlje predhodno ustvarjen osnutek sporočila. \\
\hline

\multicolumn{2}{|l|}{\textbf{Mape poštnega predala}} \\
\hline

\texttt{GET /me/mailFolders} &
Pridobi seznam map poštnega predala. \\
\hline

\texttt{POST /me/mailFolders} &
Ustvari novo mapo poštnega predala. \\
\hline

\texttt{GET ../\{mailFolder-id\}} &
Pridobi podatke o določeni mapi. \\
\hline

\texttt{PATCH ../\{mailFolder-id\}} &
Posodobi lastnosti mape. \\
\hline

\texttt{DELETE ../\{mailFolder-id\}} &
Izbriše mapo poštnega predala. \\
\hline

\multicolumn{2}{|l|}{\textbf{Sporočila v mapi}} \\
\hline

\texttt{GET ../\{mailFolder-id\}/messages} &
Pridobi sporočila iz izbrane mape. \\
\hline

\texttt{POST ../\{mailFolder-id\}/messages} &
Ustvari sporočilo v izbrani mapi. \\
\hline

\multicolumn{2}{|l|}{\textbf{Premikanje in kopiranje}} \\
\hline

\texttt{POST /me/messages/\{message-id\}
/move} &
Premakne sporočilo v drugo mapo. \\
\hline

\texttt{POST /me/messages/\{message-id\}
/copy} &
Kopira sporočilo v drugo mapo. \\
\hline

\multicolumn{2}{|l|}{\textbf{Priponke}} \\
\hline

\texttt{GET ../attachments} &
Pridobi priponke sporočila. \\
\hline

\texttt{POST ../attachments} &
Doda priponko sporočilu. \\
\hline

\texttt{GET ../attachments/\{attachment-id\}} &
Pridobi določeno priponko. \\
\hline

\texttt{DELETE ../attachments/\{attachment-id\}} &
Izbriše priponko. \\
\hline

\multicolumn{2}{|l|}{\textbf{Osnutki}} \\
\hline

\texttt{POST ../createReply} &
Ustvari osnutek odgovora na sporočilo. \\
\hline

\texttt{POST ../createReplyAll} &
Ustvari osnutek odgovora vsem prejemnikom. \\
\hline

\texttt{POST ../createForward} &
Ustvari osnutek za posredovanje sporočila. \\
\hline

\multicolumn{2}{|l|}{\textbf{Odgovarjanje in posredovanje}} \\
\hline

\texttt{POST ../reply} &
Odgovori pošiljatelju sporočila. \\
\hline

\texttt{POST ../replyAll} &
Odgovori vsem prejemnikom sporočila. \\
\hline

\texttt{POST ../forward} &
Posreduje sporočilo drugim prejemnikom. \\
\hline

\multicolumn{2}{|l|}{\textbf{Posebne operacije}} \\
\hline

\texttt{POST ../markAsJunk} &
Označi sporočilo kot nezaželeno pošto. \\
\hline

\texttt{POST ../unmarkAsJunk} &
Odstrani oznako nezaželene pošte. \\
\hline

\end{longtable}


\subsection{Koledarji in dogodki}

Koledarski del Microsoft Graph temelji na virih \texttt{Calendar}, \texttt{Event} in \texttt{CalendarView}. API omogoča popoln dostop do Outlook koledarjev ter upravljanje dogodkov, sestankov in rezervacij virov.

Dogodki vsebujejo številne atribute, med katerimi so čas začetka in konca, seznam udeležencev, lokacija, opis, ponavljajoči se vzorci ter informacije o spletnih srečanjih. Microsoft Graph podpira ustvarjanje dogodkov z integracijo Microsoft Teams, pri čemer se samodejno ustvari povezava za videokonferenčni sestanek.

Posebej uporabna funkcionalnost je metoda \texttt{findMeetingTimes}, ki analizira razpoložljivost vseh udeležencev in predlaga optimalne termine za sestanek. Pri tem sistem upošteva delovni čas uporabnikov, časovne pasove, obstoječe rezervacije in druge omejitve.

Za spremljanje sprememb koledarjev je mogoče uporabiti naročnine (angl. Subscriptions) in spletne kljuke (angl. Webhooks), ki aplikacijo obvestijo ob ustvarjanju, spremembi ali brisanju dogodkov.

\begin{longtable}{|p{7cm}|p{7.5cm}|}
\caption{Microsoft Graph API endpointi za delo s koledarji}
\label{tab:calendar-endpoints} \\
\hline
\textbf{Endpoint} & \textbf{Opis} \\
\hline
\endfirsthead

\multicolumn{2}{c}%
{{\tablename\ \thetable{} -- nadaljevanje}} \\
\hline
\textbf{Endpoint} & \textbf{Opis} \\
\hline
\endhead

\hline
\endlastfoot

\multicolumn{2}{|l|}{\textbf{Koledarji}} \\
\hline

\texttt{GET /me/calendar} &
Pridobi privzeti koledar prijavljenega uporabnika. \\
\hline

\texttt{GET /users/\{user-id\}/calendar} &
Pridobi privzeti koledar uporabnika. \\
\hline

\texttt{GET /groups/\{group-id\}/calendar} &
Pridobi privzeti koledar skupine Microsoft 365. \\
\hline

\texttt{GET /me/calendars} &
Pridobi seznam uporabnikovih koledarjev. \\
\hline

\texttt{POST /me/calendars} &
Ustvari nov koledar. \\
\hline

\texttt{GET /me/calendars/\{calendar-id\}} &
Pridobi določen koledar. \\
\hline

\texttt{PATCH /me/calendars/\{calendar-id\}} &
Posodobi lastnosti koledarja. \\
\hline

\texttt{DELETE /me/calendars/\{calendar-id\}} &
Izbriše koledar. \\
\hline

\multicolumn{2}{|l|}{\textbf{Skupine koledarjev}} \\
\hline

\texttt{GET /me/calendarGroups} &
Pridobi skupine koledarjev uporabnika. \\
\hline

\texttt{POST /me/calendarGroups} &
Ustvari novo skupino koledarjev. \\
\hline

\texttt{GET /me/calendarGroups/\{group-id\}} &
Pridobi določeno skupino koledarjev. \\
\hline

\texttt{PATCH ../\{group-id\}} &
Posodobi skupino koledarjev. \\
\hline

\texttt{DELETE ../\{group-id\}} &
Izbriše skupino koledarjev. \\
\hline

\texttt{GET ../\{group-id\}/calendars} &
Pridobi koledarje v skupini. \\
\hline

\texttt{POST ../\{group-id\}/calendars} &
Ustvari koledar v skupini. \\
\hline

\multicolumn{2}{|l|}{\textbf{Dogodki}} \\
\hline

\texttt{GET /me/events} &
Pridobi vse dogodke uporabnika. \\
\hline

\texttt{POST /me/events} &
Ustvari nov dogodek. \\
\hline

\texttt{GET /me/events/\{event-id\}} &
Pridobi določen dogodek. \\
\hline

\texttt{PATCH /me/events/\{event-id\}} &
Posodobi dogodek. \\
\hline

\texttt{DELETE /me/events/\{event-id\}} &
Izbriše dogodek. \\
\hline

\texttt{GET /me/calendar/events} &
Pridobi dogodke privzetega koledarja. \\
\hline

\texttt{POST /me/calendar/events} &
Ustvari dogodek v privzetem koledarju. \\
\hline

\texttt{GET /me/calendars/\{calendar-id\}
/events} &
Pridobi dogodke določenega koledarja. \\
\hline

\texttt{POST /me/calendars/\{calendar-id\}
/events} &
Ustvari dogodek v določenem koledarju. \\
\hline

\multicolumn{2}{|l|}{\textbf{Pogled koledarja}} \\
\hline

\texttt{GET /me/calendarView
?startDateTime=\{start\}\&endDateTime=\{end\}} &
Pridobi dogodke v podanem časovnem obdobju. \\
\hline

\texttt{GET /me/calendars
/\{calendar-id\}/calendarView
?startDateTime=\{start\}\&endDateTime=\{end\}} &
Pridobi pogled določenega koledarja za izbrano obdobje. \\
\hline

\multicolumn{2}{|l|}{\textbf{Ponavljajoči dogodki}} \\
\hline

\texttt{GET /me/events/\{event-id\}
/instances
?startDateTime=\{start\}\&endDateTime=\{end\}} &
Pridobi posamezne pojavitve ponavljajočega dogodka. \\
\hline

\multicolumn{2}{|l|}{\textbf{Odgovori na povabila}} \\
\hline

\texttt{POST /me/events/\{event-id\}/accept} &
Sprejme povabilo na dogodek. \\
\hline

\texttt{POST /me/events/\{event-id\}
/tentativelyAccept} &
Dogodek sprejme pogojno. \\
\hline

\texttt{POST /me/events/\{event-id\}/decline} &
Zavrne povabilo na dogodek. \\
\hline

\texttt{POST /me/events/\{event-id\}/cancel} &
Prekliče organiziran dogodek. \\
\hline

\texttt{POST /me/events/\{event-id\}/forward} &
Posreduje povabilo drugim uporabnikom. \\
\hline

\multicolumn{2}{|l|}{\textbf{Priponke dogodkov}} \\
\hline

\texttt{GET /me/events/\{event-id\}
/attachments} &
Pridobi priponke dogodka. \\
\hline

\texttt{POST /me/events/\{event-id\}
/attachments} &
Doda priponko dogodku. \\
\hline

\texttt{GET /me/events/\{event-id\}
/attachments/\{attachment-id\}} &
Pridobi določeno priponko. \\
\hline

\texttt{DELETE /me/events/\{event-id\}
/attachments/\{attachment-id\}} &
Izbriše priponko. \\
\hline

\multicolumn{2}{|l|}{\textbf{Dovoljenja koledarja}} \\
\hline

\texttt{GET /me/calendar
/calendarPermissions} &
Pridobi dovoljenja privzetega koledarja. \\
\hline

\texttt{GET /me/calendars/\{calendar-id\}
/calendarPermissions} &
Pridobi dovoljenja izbranega koledarja. \\
\hline

\texttt{POST /me/calendars/\{calendar-id\}
/calendarPermissions} &
Doda novo dovoljenje za koledar. \\
\hline

\texttt{PATCH /me/calendars/\{calendar-id\}
/calendarPermissions/\{permission-id\}} &
Posodobi dovoljenje. \\
\hline

\texttt{DELETE /me/calendars/\{calendar-id\}
/calendarPermissions/\{permission-id\}} &
Odstrani dovoljenje. \\
\hline

\end{longtable}

\subsection{Datoteke in dokumenti}

Microsoft Graph omogoča enoten dostop do dokumentov v storitvah OneDrive for Business in SharePoint Online preko virov \texttt{Drive} in \texttt{DriveItem}. Objekt \texttt{Drive} predstavlja logični dokumentni repozitorij, medtem ko \texttt{DriveItem} predstavlja posamezno datoteko ali mapo.

Datoteke je mogoče pridobivati preko različnih identifikatorjev, poti ali URL naslovov. Primer dostopa do vsebine datoteke:

\begin{verbatim}
/drives/{drive-id}/items/{item-id}/content
\end{verbatim}

Microsoft Graph podpira nalaganje datotek različnih velikosti. Pri manjših datotekah se uporablja enostaven prenos preko ene HTTP zahteve, medtem ko se za večje datoteke uporablja mehanizem nalaganja po segmentih (angl. Upload Session), ki omogoča nadaljevanje prenosa v primeru prekinitve povezave.

Pomembna funkcionalnost je upravljanje dovoljenj in povezav za skupno rabo. Aplikacije lahko ustvarjajo časovno omejene povezave, določajo pravice za ogled ali urejanje dokumentov ter spremljajo dostop do datotek. Poleg tega Microsoft Graph omogoča spremljanje sprememb preko mehanizma delta poizvedb (angl. Delta Query), kar je posebej uporabno pri sinhronizaciji dokumentov med sistemi.

\begin{longtable}{|p{5.5cm}|p{8.5cm}|}
\caption{Ključni Microsoft Graph API endpointi za delo z dokumenti in datotekami}
\label{tab:file-document-endpoints} \\
\hline
\textbf{Endpoint} & \textbf{Opis} \\
\hline
\endfirsthead

\multicolumn{2}{c}%
{{\tablename\ \thetable{} -- nadaljevanje}} \\
\hline
\textbf{Endpoint} & \textbf{Opis} \\
\hline
\endhead

\hline
\endlastfoot

\texttt{GET /me/drive} &
Pridobi OneDrive trenutnega uporabnika. \\
\hline

\texttt{GET /drives} &
Pridobi seznam pogonov, do katerih ima uporabnik dostop. \\
\hline

\texttt{GET /drives/\{drive-id\}} &
Pridobi informacije o določenem pogonu. \\
\hline

\texttt{GET /me/drive/root} &
Pridobi korensko mapo uporabnikovega OneDrive-a. \\
\hline

\texttt{GET /me/drive/root/children} &
Pridobi seznam datotek in map v korenski mapi. \\
\hline

\texttt{GET /drives/\{drive-id\}/items/
\{item-id\}} &
Pridobi informacije o določeni datoteki ali mapi. \\
\hline

\texttt{GET /drives/\{drive-id\}/items/
\{item-id\}/children} &
Pridobi vsebino določene mape. \\
\hline

\texttt{PUT /drives/\{drive-id\}/items/
\{parent-id\}:/\{filename\}:/content} &
Naloži novo datoteko v izbrano mapo ali zamenja obstoječo datoteko. \\
\hline

\texttt{GET ../\{item-id\}/content} &
Prenese vsebino datoteke. \\
\hline

\texttt{PATCH ..\{item-id\}} &
Posodobi lastnosti datoteke ali mape (npr. ime). \\
\hline

\texttt{DELETE ../\{item-id\}} &
Izbriše datoteko ali mapo. \\
\hline

\texttt{POST ../\{item-id\}/copy} &
Ustvari kopijo datoteke ali mape. \\
\hline

\texttt{POST ../\{item-id\}/move} &
Premakne datoteko ali mapo na drugo mesto znotraj pogona. \\
\hline

\texttt{POST ../\{item-id\}/createLink} &
Ustvari povezavo za skupno rabo datoteke ali mape. \\
\hline

\texttt{POST ../\{item-id\}/invite} &
Podeli dostop do datoteke ali mape drugim uporabnikom. \\
\hline

\texttt{GET /shares/\{share-id\}
/driveItem} &
Pridobi datoteko ali mapo prek povezave za skupno rabo. \\
\hline

\texttt{GET /sites/\{site-id\}/drive} &
Pridobi dokumentno knjižnico SharePoint mesta. \\
\hline

\texttt{GET /sites/\{site-id\}/drive
/root/children} &
Pridobi dokumente in mape v korenski mapi SharePoint knjižnice. \\
\hline

\texttt{GET /groups/\{group-id\}/drive} &
Pridobi dokumentno knjižnico skupine Microsoft 365. \\
\hline

\texttt{GET /groups/\{group-id\}/drive
/root/children} &
Pridobi datoteke in mape v dokumentni knjižnici skupine. \\
\hline

\texttt{GET /drives/\{drive-id\}/items
/\{item-id\}/permissions} &
Pridobi dovoljenja za datoteko ali mapo. \\
\hline

\texttt{POST ../permissions} &
Ustvari novo dovoljenje za dostop do datoteke ali mape. \\
\hline

\texttt{DELETE ../permissions/\{perm-id\}} &
Odstrani določeno dovoljenje za datoteko ali mapo. \\
\hline

\end{longtable}

\subsection{Excel in poslovni podatki}

Microsoft Graph vključuje namenski Excel API, ki omogoča neposredno delo z delovnimi zvezki brez uporabe namizne aplikacije Excel. Delovni zvezek je predstavljen kot hierarhična struktura objektov, ki vključuje delovne liste, tabele, obsege celic, grafikone in vrtilne tabele.

Dostop do podatkov poteka preko virov, kot so:

\begin{verbatim}
/workbook/worksheets
/workbook/tables
/workbook/range
\end{verbatim}

API omogoča branje in zapisovanje vrednosti celic, izvajanje formul, osveževanje podatkovnih povezav ter upravljanje struktur tabel.
Posebna prednost je možnost izvajanja operacij neposredno na strežniški strani, kar zmanjšuje potrebo po prenosu celotnih datotek.

Za zagotavljanje konsistentnosti podatkov Microsoft Graph podpira koncept sej (angl. Workbook Session), ki omogoča izvajanje več operacij nad delovnim zvezkom znotraj istega konteksta. Takšen pristop izboljša zmogljivost in preprečuje konflikte pri sočasnem dostopu več uporabnikov.

Excel API je posebej uporaben pri avtomatizaciji poročanja, pripravi finančnih analiz, integraciji poslovnih podatkov in gradnji analitičnih rešitev, kjer Excel predstavlja podatkovni vir ali uporabniški vmesnik za poslovne uporabnike.

\begin{longtable}{|p{7cm}|p{7.5cm}|}
\caption{Ključni Microsoft Graph API endpointi za delo z Excelovimi dokumenti}
\label{tab:excel-endpoints} \\
\hline
\textbf{Endpoint} & \textbf{Opis} \\
\hline
\endfirsthead

\multicolumn{2}{c}%
{{\tablename\ \thetable{} -- nadaljevanje}} \\
\hline
\textbf{Endpoint} & \textbf{Opis} \\
\hline
\endhead

\hline
\endlastfoot

\multicolumn{2}{|l|}{\textbf{Delovni zvezek}} \\
\hline

\texttt{GET /me/drive/items/\{item-id\}
/workbook} &
Pridobi Excelov delovni zvezek uporabnika. \\
\hline

\texttt{GET /drives/\{drive-id\}/items
/\{item-id\}/workbook} &
Pridobi Excelov delovni zvezek iz določenega pogona. \\
\hline

\texttt{POST ../workbook/createSession} &
Ustvari sejo za delo z Excelovim dokumentom. \\
\hline

\texttt{POST ../workbook/closeSession} &
Zapre aktivno Excelovo sejo. \\
\hline

\multicolumn{2}{|l|}{\textbf{Delovni listi}} \\
\hline

\texttt{GET .drives/\{drive-id\}/items
/\{item-id\}/workbook/worksheets} &
Pridobi seznam delovnih listov v dokumentu. \\
\hline

\texttt{GET ../\{worksheet-id\}} &
Pridobi podatke o določenem delovnem listu. \\
\hline

\texttt{POST ../worksheets/add} &
Ustvari nov delovni list. \\
\hline

\texttt{PATCH ../worksheets/\{worksheet-id\}} &
Posodobi lastnosti delovnega lista. \\
\hline

\texttt{DELETE ../worksheets/\{worksheet-id\}} &
Izbriše delovni list. \\
\hline

\multicolumn{2}{|l|}{\textbf{Območja celic}} \\
\hline

\texttt{GET ../range(address='A1:C10')} &
Prebere vrednosti, formule in oblikovanje območja celic. \\
\hline

\texttt{PATCH ../range(address='A1:C10')} &
Posodobi vrednosti ali formule v območju celic. \\
\hline

\texttt{POST ../range(address='A1:C10')
/clear} &
Počisti vsebino ali oblikovanje izbranega območja. \\
\hline

\texttt{POST ../range(address='A1:C10')
/insert} &
Vstavi nove celice v območje. \\
\hline

\texttt{POST ../range(address='A1:C10')
/delete} &
Odstrani celice iz območja. \\
\hline

\multicolumn{2}{|l|}{\textbf{Tabele}} \\
\hline

\texttt{GET ../workbook/tables} &
Pridobi seznam tabel v delovnem zvezku. \\
\hline

\texttt{POST ../tables/add} &
Ustvari novo tabelo iz območja celic. \\
\hline

\texttt{GET ../tables/\{table-id\}} &
Pridobi podatke o določeni tabeli. \\
\hline

\texttt{PATCH ../tables/\{table-id\}} &
Posodobi lastnosti tabele. \\
\hline

\texttt{DELETE ../tables/\{table-id\}} &
Izbriše tabelo. \\
\hline

\multicolumn{2}{|l|}{\textbf{Stolpci in vrstice tabel}} \\
\hline

\texttt{GET ../tables/\{table-id\}/columns} &
Pridobi stolpce izbrane tabele. \\
\hline

\texttt{POST ../\{table-id\}/columns/add} &
Doda nov stolpec v tabelo. \\
\hline

\texttt{GET ../\{table-id\}/rows} &
Pridobi vrstice Excelove tabele. \\
\hline

\texttt{POST ../\{table-id\}/rows/add} &
Doda novo vrstico v tabelo. \\
\hline

\texttt{DELETE ../\{table-id\}/rows/\{index\}} &
Izbriše določeno vrstico tabele. \\
\hline

\multicolumn{2}{|l|}{\textbf{Imenovani obsegi}} \\
\hline

\texttt{GET ../workbook/names} &
Pridobi imenovane obsege v dokumentu. \\
\hline

\texttt{POST ../workbook/names/add} &
Ustvari nov imenovan obseg. \\
\hline

\multicolumn{2}{|l|}{\textbf{Grafi}} \\
\hline

\texttt{GET ../workbook/charts} &
Pridobi seznam grafov. \\
\hline

\texttt{POST ../workbook/worksheets
/\{worksheet-id\}/charts/add} &
Ustvari nov graf. \\
\hline

\texttt{GET ../charts/\{chart-id\}} &
Pridobi podatke o določenem grafu. \\
\hline

\texttt{PATCH ../charts/\{chart-id\}} &
Posodobi lastnosti grafa. \\
\hline

\texttt{DELETE ../charts/\{chart-id\}} &
Izbriše graf. \\
\hline

\texttt{GET ../range(address='A1:C10')
/format} &
Pridobi oblikovanje območja celic. \\
\hline

\texttt{PATCH ../range(address='A1:C10')
/format} &
Posodobi oblikovanje območja celic. \\
\hline

\texttt{POST ../range(address='A1:C10')
/sort} &
Razvrsti podatke v izbranem območju. \\
\hline

\texttt{POST ../range(address='A1:C10')
/filter} &
Uporabi filter nad podatki. \\
\hline

\texttt{POST ../functions/\{function-name\}} &
Izvede vgrajeno Excelovo funkcijo. \\
\hline

\end{longtable}

\subsection{OneNote}

OneNote API omogoča dostop do beležnic, odsekov in posameznih strani znotraj storitve Microsoft OneNote. Podatkovni model temelji na hierarhiji objektov \texttt{Notebook}, \texttt{Section} in \texttt{Page}, kar omogoča strukturirano organizacijo vsebin.

Posebnost OneNote API-ja je uporaba HTML vsebine za predstavitev strani. Pri ustvarjanju ali posodabljanju zapiskov aplikacija pošilja HTML dokumente, ki lahko vsebujejo oblikovano besedilo, slike, povezave, sezname in druge elemente. Takšen pristop omogoča bogato predstavitev vsebine brez potrebe po uporabi lastnega formata dokumentov.

Microsoft Graph omogoča tudi iskanje po vsebini beležnic, pridobivanje metapodatkov ter upravljanje prilog. Zaradi teh lastnosti je OneNote primeren za implementacijo sistemov za upravljanje znanja, projektne dokumentacije in zapisnikov sestankov.

\begin{longtable}{|p{6cm}|p{8cm}|}
\caption{Ključni Microsoft Graph API endpointi za delo z OneNote dokumenti}
\label{tab:onenote-endpoints} \\
\hline
\textbf{Endpoint} & \textbf{Opis} \\
\hline
\endfirsthead

\multicolumn{2}{c}%
{{\tablename\ \thetable{} -- nadaljevanje}} \\
\hline
\textbf{Endpoint} & \textbf{Opis} \\
\hline
\endhead

\hline
\endlastfoot

\multicolumn{2}{|l|}{\textbf{Upravljanje zvezkov}} \\
\hline

\texttt{GET /me/onenote/notebooks} &
Pridobi seznam zvezkov OneNote trenutnega uporabnika. \\
\hline

\texttt{POST /me/onenote/notebooks} &
Ustvari nov OneNote zvezek. \\
\hline

\texttt{GET ../\{notebook-id\}} &
Pridobi podatke o določenem OneNote zvezku. \\
\hline

\texttt{PATCH ../\{notebook-id\}} &
Posodobi lastnosti OneNote zvezka. \\
\hline

\texttt{DELETE ../\{notebook-id\}} &
Izbriše OneNote zvezek. \\
\hline

\texttt{GET ../\{notebook-id\}/sections} &
Pridobi razdelke izbranega zvezka. \\
\hline

\texttt{POST ../\{notebook-id\}/sections} &
Ustvari nov razdelek v zvezku. \\
\hline

\multicolumn{2}{|l|}{\textbf{Organizacija vsebine}} \\
\hline

\texttt{GET /me/onenote/sections} &
Pridobi vse razdelke OneNote uporabnika. \\
\hline

\texttt{GET ../\{section-id\}} &
Pridobi podatke o določenem razdelku. \\
\hline

\texttt{PATCH ../\{section-id\}} &
Posodobi lastnosti razdelka. \\
\hline

\texttt{DELETE ../\{section-id\}} &
Izbriše razdelek. \\
\hline

\texttt{GET ../\{section-id\}/pages} &
Pridobi strani v izbranem razdelku. \\
\hline

\texttt{POST ../\{section-id\}/pages} &
Ustvari novo stran v razdelku. \\
\hline

\multicolumn{2}{|l|}{\textbf{Branje in zapisovanje zapiskov}} \\
\hline

\texttt{GET /me/onenote/pages} &
Pridobi vse OneNote strani uporabnika. \\
\hline

\texttt{GET ../\{page-id\}} &
Pridobi podatke o določeni strani. \\
\hline

\texttt{PATCH ../\{page-id\}} &
Posodobi lastnosti strani. \\
\hline

\texttt{DELETE ../\{page-id\}} &
Izbriše stran. \\
\hline

\texttt{GET ../\{page-id\}/content} &
Pridobi HTML vsebino strani OneNote. \\
\hline

\texttt{PATCH ../\{page-id\}/content} &
Posodobi vsebino strani z uporabo ukazov HTML. \\
\hline

\texttt{GET ../\{page-id\}/preview} &
Pridobi predogled strani. \\
\hline

\multicolumn{2}{|l|}{\textbf{Slike in priponke}} \\
\hline

\texttt{GET /me/onenote/resources} &
Pridobi vire, povezane z OneNote stranmi (slike, datoteke in drugi objekti). \\
\hline

\texttt{GET ../\{resource-id\}} &
Pridobi določen vir strani. \\
\hline

\texttt{GET ../\{resource-id\}/\$value} &
Prenese binarno vsebino vira (npr. sliko ali datoteko). \\
\hline

\multicolumn{2}{|l|}{\textbf{Skupine razdelkov}} \\
\hline

\texttt{GET /me/onenote/notebooks
/\{notebook-id\}/sectionGroups} &
Pridobi skupine razdelkov v zvezku. \\
\hline

\texttt{POST /me/onenote/notebooks
/\{notebook-id\}/sectionGroups} &
Ustvari novo skupino razdelkov. \\
\hline

\texttt{GET /me/onenote/sectionGroups
/\{sectionGroup-id\}} &
Pridobi podatke o skupini razdelkov. \\
\hline

\texttt{PATCH /me/onenote/sectionGroups
/\{sectionGroup-id\}} &
Posodobi skupino razdelkov. \\
\hline

\texttt{DELETE /me/onenote/sectionGroups
/\{sectionGroup-id\}} &
Izbriše skupino razdelkov. \\
\hline

\texttt{GET /me/onenote/sectionGroups
/\{sectionGroup-id\}/sections} &
Pridobi razdelke znotraj skupine razdelkov. \\
\hline

\multicolumn{2}{|l|}{\textbf{Dostop do zvezkov}} \\
\hline

\texttt{GET /users/\{user-id\}
/onenote/notebooks} &
Pridobi OneNote zvezke določenega uporabnika. \\
\hline

\texttt{GET /groups/\{group-id\}
/onenote/notebooks} &
Pridobi OneNote zvezke skupine Microsoft 365. \\
\hline

\texttt{GET /sites/\{site-id\}
/onenote/notebooks} &
Pridobi OneNote zvezke, povezane s SharePoint mestom. \\
\hline

\end{longtable}

\subsection{Microsoft Teams}

Microsoft Teams predstavlja osrednjo platformo za sodelovanje in komunikacijo v okolju Microsoft 365. Microsoft Graph omogoča dostop do vseh ključnih komponent sistema Teams, vključno z ekipami, kanali, člani, sporočili, datotekami in spletnimi sestanki.

Podatkovni model temelji na virih \texttt{Team}, \texttt{Channel}, \texttt{Chat} in \texttt{OnlineMeeting}. Ker je vsaka ekipa povezana z Microsoft 365 skupino, so vsi povezani viri dostopni preko skupnega identitetnega modela.

Microsoft Graph omogoča ustvarjanje ekip in kanalov, upravljanje članstva ter dostop do pogovorov. Za avtomatizacijo poslovnih procesov je posebej pomembna možnost pošiljanja sporočil v kanale in klepete ter integracija z aplikacijami in roboti.

Pri spletnih sestankih API omogoča ustvarjanje, načrtovanje in spremljanje srečanj Teams. Razvijalci lahko pridobivajo informacije o udeležencih, času trajanja sestankov in drugih metapodatkih, kar omogoča razvoj naprednih rešitev za upravljanje komunikacije in sodelovanja.

Zaradi tesne povezanosti s SharePointom, Outlookom, OneDriveom in Microsoft Entra ID predstavlja Teams pomembno integracijsko točko za razvoj poslovnih aplikacij, ki združujejo komunikacijo, dokumente in poslovne procese v enotnem uporabniškem okolju.

\begin{longtable}{|p{6cm}|p{8cm}|}
\caption{Ključni Microsoft Graph API endpointi za delo z Microsoft Teams}
\label{tab:teams-endpoints} \\
\textbf{Endpoint} & \textbf{Opis} \\
\hline
\endfirsthead

\multicolumn{2}{c}%
{{\tablename\ \thetable{} -- nadaljevanje}} \\
\hline
\textbf{Endpoint} & \textbf{Opis} \\
\hline
\endhead

\multicolumn{2}{|l|}{\textbf{Ustvarjanje in upravljanje ekip}} \\
\hline

\texttt{GET /teams} &
Pridobi seznam ekip Microsoft Teams, do katerih ima uporabnik dostop. \\
\hline

\texttt{GET /teams/\{team-id\}} &
Pridobi podatke o določeni ekipi. \\
\hline

\texttt{POST /teams} &
Ustvari novo ekipo Microsoft Teams. \\
\hline

\texttt{PATCH /teams/\{team-id\}} &
Posodobi lastnosti ekipe. \\
\hline

\texttt{DELETE /teams/\{team-id\}} &
Izbriše ekipo. \\
\hline

\texttt{GET /groups/\{group-id\}/team} &
Pridobi ekipo, povezano z Microsoft 365 skupino. \\
\hline

\texttt{PUT /groups/\{group-id\}/team} &
Ustvari ekipo iz obstoječe Microsoft 365 skupine. \\
\hline

\multicolumn{2}{|l|}{\textbf{Kanali in zasebni kanali}} \\
\hline

\texttt{GET /teams/\{team-id\}/channels} &
Pridobi seznam kanalov v ekipi. \\
\hline

\texttt{POST /teams/\{team-id\}/channels} &
Ustvari nov kanal v ekipi. \\
\hline

\texttt{GET ../\{channel-id\}} &
Pridobi podatke o določenem kanalu. \\
\hline

\texttt{PATCH ../\{channel-id\}} &
Posodobi podatke kanala. \\
\hline

\texttt{DELETE ../\{channel-id\}} &
Izbriše kanal. \\
\hline

\multicolumn{2}{|l|}{\textbf{Člani ekip in kanalov}} \\
\hline

\texttt{GET /teams/\{team-id\}
/\{channel-id\}/members} &
Pridobi člane zasebnega kanala. \\
\hline

\texttt{POST ../members} &
Doda člana v zasebni kanal. \\
\hline

\texttt{DELETE ../members/\{membership-id\}} &
Odstrani člana iz zasebnega kanala. \\
\hline

\texttt{GET /teams/\{team-id\}/members} &
Pridobi člane ekipe. \\
\hline

\texttt{POST /teams/\{team-id\}/members} &
Doda člana v ekipo. \\
\hline

\texttt{DELETE /teams/\{team-id\}
/members/\{membership-id\}} &
Odstrani člana iz ekipe. \\
\hline

\multicolumn{2}{|l|}{\textbf{Sporočila in odgovori}} \\
\hline

\texttt{GET /teams/\{team-id\}
/\{channel-id\}/messages} &
Pridobi sporočila v kanalu. \\
\hline

\texttt{POST ../messages} &
Ustvari novo sporočilo v kanalu. \\
\hline

\texttt{GET ../messages/\{message-id\}} &
Pridobi določeno sporočilo kanala. \\
\hline

\texttt{PATCH ../messages/\{message-id\}} &
Posodobi sporočilo kanala. \\
\hline

\texttt{DELETE ../messages/\{message-id\}} &
Izbriše sporočilo kanala. \\
\hline

\texttt{POST ../messages/\{message-id\}/replies} &
Ustvari odgovor na sporočilo. \\
\hline

\multicolumn{2}{|l|}{\textbf{Zasebni klepeti}} \\
\hline

\texttt{GET /chats} &
Pridobi seznam klepetov uporabnika. \\
\hline

\texttt{POST /chats} &
Ustvari nov klepet. \\
\hline

\texttt{GET /chats/\{chat-id\}} &
Pridobi določen klepet. \\
\hline

\texttt{DELETE /chats/\{chat-id\}} &
Izbriše klepet. \\
\hline

\texttt{GET /chats/\{chat-id\}/messages} &
Pridobi sporočila klepeta. \\
\hline

\texttt{POST /chats/\{chat-id\}/messages} &
Pošlje novo sporočilo v klepet. \\
\hline

\texttt{GET /chats/\{chat-id\}/members} &
Pridobi člane klepeta. \\
\hline

\texttt{POST /chats/\{chat-id\}/members} &
Doda člana v klepet. \\
\hline

\multicolumn{2}{|l|}{\textbf{Aplikacije Teams}} \\
\hline

\texttt{GET /teams/\{team-id\}/installedApps} &
Pridobi nameščene aplikacije v ekipi. \\
\hline

\texttt{POST ../installedApps} &
Namesti aplikacijo v ekipo. \\
\hline

\texttt{DELETE ../installedApps/\{app-installation-id\}} &
Odstrani aplikacijo iz ekipe. \\
\hline

\multicolumn{2}{|l|}{\textbf{Zavihki v kanalih}} \\
\hline

\texttt{GET /teams/\{team-id\}
/channels/\{channel-id\}/tabs} &
Pridobi zavihke kanala. \\
\hline

\texttt{POST ../tabs} &
Ustvari nov zavihek v kanalu. \\
\hline

\texttt{GET ../tabs/\{tab-id\}} &
Pridobi določen zavihek. \\
\hline

\texttt{PATCH ../tabs/\{tab-id\}} &
Posodobi zavihek. \\
\hline

\texttt{DELETE ../tabs/\{tab-id\}} &
Odstrani zavihek. \\
\hline

\multicolumn{2}{|l|}{\textbf{Podatki o uporabniku}} \\
\hline

\texttt{GET /users/\{user-id\}/joinedTeams} &
Pridobi ekipe, katerih član je uporabnik. \\
\hline

\texttt{GET /me/joinedTeams} &
Pridobi ekipe trenutno prijavljenega uporabnika. \\
\hline

\texttt{GET /me/chats} &
Pridobi klepete trenutnega uporabnika. \\
\hline

\multicolumn{2}{|l|}{\textbf{Spletni sestanki}} \\
\hline

\texttt{GET /me/onlineMeetings} &
Pridobi spletne sestanke uporabnika. \\
\hline

\texttt{POST /me/onlineMeetings} &
Ustvari nov spletni sestanek Teams. \\
\hline

\texttt{GET /me/onlineMeetings/\{meeting-id\}} &
Pridobi podatke o spletnem sestanku. \\
\hline

\texttt{PATCH /me/onlineMeetings/\{meeting-id\}} &
Posodobi spletni sestanek. \\
\hline

\texttt{DELETE /me/onlineMeetings/\{meeting-id\}} &
Izbriše spletni sestanek. \\
\hline

\multicolumn{2}{|l|}{\textbf{Organizacijski podatki}} \\
\hline

\texttt{GET /communications/callRecords} &
Pridobi zapise klicev Teams. \\
\hline

\texttt{GET /communications/onlineMeetings} &
Dostop do spletnih sestankov v organizaciji. \\
\hline

\end{longtable}

\chapter{Implementacija aplikacije}

Za praktični del diplomske naloge je bila razvita demonstracijska enostranska spletna aplikacija (angl. Single Page Application – SPA), katere namen je prikaz uporabe Microsoft Graph API pri razvoju sodobnih informacijskih sistemov. Aplikacija uporabnikom omogoča prijavo z Microsoft računom, dostop do izbranih storitev okolja Microsoft 365 ter delo z dokumenti in elektronsko pošto. Pri implementaciji je bil Microsoft Graph uporabljen kot osrednji integracijski mehanizem za komunikacijo z oblačnimi storitvami Microsoft 365.

\section{Arhitektura rešitve}

Pri načrtovanju aplikacije je bila izbrana arhitektura SPA, ki omogoča izvajanje večine logike na odjemalski strani in neposredno komunikacijo z Microsoft Graph API. Takšen pristop zmanjšuje potrebo po dodatnem strežniškem sloju, poenostavlja implementacijo in omogoča hitrejši razvoj prototipa.

Uporabniški vmesnik je bil razvit z uporabo knjižnice React in programskega jezika TypeScript, razvojno okolje pa temelji na orodju Vite. Za avtentikacijo uporabnikov in pridobivanje dostopnih žetonov je bila uporabljena knjižnica Microsoft Authentication Library (MSAL), ki zagotavlja podporo za integracijo z Microsoft Entra ID.

Arhitektura rešitve temelji na neposredni komunikaciji med odjemalsko aplikacijo in storitvami Microsoft 365 preko Microsoft Graph API. Uporabnik se najprej avtenticira v sistem Microsoft Entra ID, nato pa aplikacija pridobi dostopni žeton, ki se uporablja za izvajanje zahtevkov do storitev Microsoft 365.

\section{Avtentikacija in avtorizacija}

Za prijavo uporabnikov je bila uporabljena identitetna platforma Microsoft Entra ID. Aplikacija je bila registrirana v administrativnem portalu Microsoft Entra, kjer so bili konfigurirani identifikatorji aplikacije, preusmeritveni naslovi in zahtevana dovoljenja.

Postopek prijave temelji na protokolu OAuth 2.0 Authorization Code Flow s podporo za mehanizem PKCE (angl. Proof Key for Code Exchange), ki predstavlja priporočeni pristop za enostranske spletne aplikacije. Takšna zasnova omogoča varno pridobivanje dostopnih žetonov brez potrebe po shranjevanju občutljivih podatkov na odjemalski strani.

Za delovanje aplikacije so bila uporabljena dovoljenja za dostop do uporabniškega profila, pošiljanje elektronske pošte ter delo z dokumenti v storitvah OneDrive in SharePoint.

\subsection{Integracija z Microsoft Graph}

Microsoft Graph predstavlja osrednji komunikacijski sloj aplikacije in omogoča dostop do različnih storitev okolja Microsoft 365 preko enotnega programskega vmesnika. 
V okviru implementacije je bil uporabljen za pridobivanje podatkov o prijavljenem uporabniku, dostop do dokumentov, ustvarjanje povezav za skupno rabo ter pošiljanje elektronske pošte.

Komunikacija z Microsoft Graph poteka preko REST zahtevkov, pri čemer se podatki izmenjujejo v formatu JSON. Takšen pristop omogoča enostavno integracijo z različnimi storitvami Microsoft 365 ter zmanjšuje kompleksnost razvoja aplikacije.

\section{Delo z dokumenti}

Ena izmed osrednjih funkcionalnosti aplikacije je delo z dokumenti, shranjenimi v storitvah OneDrive in SharePoint. Uporabnik lahko izbere predlogo dokumenta, na podlagi katere aplikacija ustvari novo kopijo dokumenta in jo shrani v uporabnikov prostor za hrambo datotek.

Posebna pozornost je bila namenjena delu z Excel dokumenti.
Microsoft Graph omogoča neposreden dostop do strukture Excelovih delovnih zvezkov, delovnih listov in celic, zato je mogoče podatke pridobivati in obdelovati brez prenosa celotne datoteke. Takšen pristop omogoča učinkovito ekstrakcijo vsebine in njeno pretvorbo v obliko, primerno za nadaljnjo obdelavo v informacijskem sistemu.

Pri dokumentih Word in PowerPoint Microsoft Graph omogoča predvsem upravljanje datotek in njihovih metapodatkov, medtem ko je za podrobnejšo obdelavo vsebine potrebno uporabiti dodatne knjižnice za delo z dokumenti Office Open XML.

\section{Pošiljanje elektronske pošte}

Aplikacija omogoča pošiljanje elektronskih sporočil preko storitve Outlook.
Po pripravi dokumenta lahko uporabnik vnese prejemnike, zadevo in vsebino sporočila, nato pa aplikacija preko Microsoft Graph izvede pošiljanje elektronske pošte.

Sporočilom je mogoče priložiti dokumente ali vključiti povezave do datotek, shranjenih v storitvah OneDrive oziroma SharePoint.
Takšna funkcionalnost omogoča učinkovito povezovanje dokumentnega sistema in elektronske komunikacije znotraj enotnega uporabniškega okolja.

\section{Ocena rešitve}

Razvita aplikacija prikazuje praktično uporabo Microsoft Graph API pri razvoju sodobnih spletnih informacijskih sistemov. Uporaba Microsoft Entra ID za upravljanje identitet, Microsoft Graph za dostop do storitev ter obstoječe infrastrukture Microsoft 365 omogoča razvoj rešitev brez potrebe po lastnih strežnikih za upravljanje uporabnikov, dokumentov ali elektronske pošte.

Rezultati implementacije potrjujejo, da Microsoft Graph predstavlja učinkovito platformo za integracijo različnih storitev Microsoft 365 ter omogoča razvoj aplikacij, ki združujejo upravljanje dokumentov, komunikacijo in uporabniške podatke v enotnem informacijskem sistemu.




\chapter{Sklepne ugotovitve}




\cleardoublepage
\addcontentsline{toc}{chapter}{Literatura}
\bibliography{literatura}
\bibliographystyle{plainnat}

\end{document}

